La región sólida $E$ está acotada por el cilindro circular $x^2 + y^2 = 4$, el plano $z = 0$ y el plano $z = 3$.
\begin{enumerate}
    \item Describir en palabras la forma del sólido $E$.
    \item Determinar si el punto $Q = (1,\, 1,\, 2)$ está dentro, fuera o sobre la frontera de $E$.
    \item Encontrar la longitud del segmento de la recta $\vec{r}(t) = (1,\, 0,\, t)$, $t \in \mathbb{R}$, que se encuentra dentro de $E$.
    \item Dos cilindros, $x^2 + y^2 = 4$ y $x^2 + z^2 = 4$, se intersectan. Encontrar y describir la curva de intersección en la región $x \geq 0$.

    \textit{Sugerencia:} En la curva de intersección se cumple $y^2 = z^2$; usar esto para parametrizar.
\end{enumerate}